\begin{abstract}
HomeWorkout is an Android application that helps people plan and complete structured exercise sessions without depending on a gym coach or a continuous internet connection. The application combines a local exercise library, profile-based plan recommendation, editable multi-day workout plans, guided workout playback, progress tracking, reminders, achievements, weight monitoring, and an experimental AI fitness assistant.

The project is implemented in Kotlin with Jetpack Compose. It follows a three-layer structure made of UI, domain, and data layers. ViewModels expose observable state, use cases contain focused application rules, repositories separate domain models from storage and network details, and Room provides the local source of truth. The bundled dataset contains 771 usable exercises and the recommendation catalog contains 16 plan templates. Most training functions continue to work offline after the first installation. The chatbot is the main network-dependent function and calls the Groq API through OkHttp.

This report explains the problem addressed by the application, its target users, the implemented user flows, architecture, database design, technology choices, setup process, and team responsibilities. It also evaluates the current implementation honestly. The main training, plan editing, settings, reminder, streak, and achievement flows are implemented. Weight and database-backed history work are available on the main development line, while the chatbot is implemented on a separate branch that has not yet been safely integrated with the latest database migration. The current prototype builds successfully, but automated test coverage and production security still require further work.
\end{abstract}
